\documentclass[11pt]{article}
\usepackage[margin=25mm]{geometry}
\usepackage{fontspec}
\setmainfont{Latin Modern Roman}
\usepackage{amsmath,amssymb,mathtools}
\usepackage{unicode-math}
\setmathfont{Latin Modern Math}
\usepackage{xurl,hyperref}
\hypersetup{colorlinks=true,urlcolor=blue}
\setcounter{secnumdepth}{0}
\setlength{\parindent}{0pt}
\setlength{\parskip}{5pt}
\setlength{\emergencystretch}{3em}
\newcommand{\sref}[2]{\hyperlink{source:#1:#2}{[S#2]}}
\newcommand{\cataloguescope}[1]{\paragraph{Recorded target.}#1}
\begin{document}
\section{Research Problems in Function Theory — Problem 7.6}\label{statement}

\subsection{Definitions and mathematical statement}

\paragraph{Shared notation.}$\mathbb N=\{1,2,\ldots\}$, and $\mathbb Z,\mathbb Q,\mathbb R,\mathbb C$ denote the integers, rationals, reals and complex numbers. Any other conventions are specified locally.


Write $\mathbb D=\{z\in\mathbb C:|z|<1\}$ and $\mathbb T=\partial\mathbb D$. Let $(a_n)_{n\geq0}$ be fixed complex coefficients with $\limsup_{n\to\infty}|a_n|^{1/n}\leq1$ and $\sum_{n=0}^{\infty}|a_n|^2=\infty$. Let $(\varepsilon_n)$ be independent random signs with $\mathbb P(\varepsilon_n=1)=\mathbb P(\varepsilon_n=-1)=1/2$, and define the holomorphic random series
\[F(z)=\sum_{n=0}^{\infty}\varepsilon_na_nz^n,\qquad z\in\mathbb D.\]
The radius assumption states explicitly that the series in the question is defined on $\mathbb D$.
\paragraph{Statement recorded in the supplied source.}
Are each of the following true: $\mathbb P(\overline{F(\mathbb D)}=\mathbb C)=1$; $\mathbb P(F(\mathbb D)=\mathbb C)=1$; and, for every fixed $w\in\mathbb C$, $\mathbb P(w\in F(\mathbb D))=1$? The source's update attributes an affirmative answer to the whole-plane assertion to Nazarov, Sodin and Nishry, which also implies the other two assertions. \sref{statement}{2}
\subsection{Short English statement}
Does every random-sign power series with nonsquare-summable coefficients almost surely have dense, or full, complex range and hit each fixed point?
\subsection{Sources}
\begin{description}
\item[\hypertarget{source:statement:1}{S1}] ULAM.AI and the UnsolvedMath Contributors, UnsolvedMath: A Curated Collection of Open Mathematics Problems, catalogue entry. CC BY 4.0. \url{https://huggingface.co/datasets/ulamai/UnsolvedMath}.
\item[\hypertarget{source:statement:2}{S2}] W. K. Hayman and E. F. Lingham, Research Problems in Function Theory (2018), Problem 7.6, complete problem and update. \url{https://arxiv.org/abs/1809.07200}.
\end{description}
\label{end:statement}
\end{document}
